\section{Structure and classification of simple Lie algebras}

\lecture{6} We begin with some linear algebra preliminaries.

\begin{definition}
A Lie algebra $\g$ is called \emph{simple} if the only ideals in $\g$ are $0$
and $\g$, and $\dim \g > 1$ (equivalently, given the first condition, $\g$ is
not abelian). A Lie algebra $\g$ is called \emph{semisimple} if it is a direct
sum of simple Lie algebras.
\end{definition}

\begin{remark}
Recall that $\C$ behaved quite differently from $\sl_2$, so it is sensible to
exclude the abelian Lie algebras.
\end{remark}

\begin{definition}
The \emph{derived algebra} of $\g$ is
$[\g, \g] = \Span\setc{[x, y]}{x, y \in \g}$.
\end{definition}

\begin{exercise}\leavevmode
\begin{enumerate}
  \item $[\g, \g]$ is an ideal;
  \item $\g/[\g, \g]$ is abelian.
\end{enumerate}
\end{exercise}

\begin{definition}
The \emph{central series} of $\g$ is given by $\g^0 = \g$,
$\g^n = [\g^{n-1}, \g]$, i.e.
\[
  \g \supseteq [\g, \g] \supseteq [[\g, \g], \g] \supseteq \cdots
\]
The \emph{derived series} of $\g$ is given by $\g^{(0)} = \g$,
$\g^{(n)} = [\g^{(n-1)}, \g^{(n-1)}]$, i.e.
\[
  \g \supseteq [\g, \g] \supseteq [[\g, \g], [\g, \g]] \supseteq \cdots
\]
\end{definition}

\begin{remark}
Clearly, $\g^{(n)} \subseteq \g^n$.
\end{remark}

\begin{definition}
$\g$ is \emph{nilpotent} if $\g^n = 0$ for some $n > 0$, i.e.\ the central
series terminates. $\g$ is \emph{solvable} if $\g^{(n)} = 0$ for some $n > 0$,
i.e.\ the derived series terminates.
\end{definition}

By the above remark, a nilpotent Lie algebra is necessarily solvable.

\begin{remark}
The term ``solvable'' comes from Galois theory (if the Galois group is
solvable, the associated polynomial is solvable by radicals). We will see
shortly where the term ``nilpotent'' comes from.
\end{remark}

\begin{example}\leavevmode
\begin{enumerate}
  \item $\fn = \Biggl\{ \begin{pmatrix}
            0 & * & \cdots & *\\
              & 0 & \cdots & *\\
              &   & \ddots & \vdots\\
              &   &        & 0
          \end{pmatrix} \in \gl_n \Biggr\}$, the strictly upper triangular
    matrices, form a nilpotent Lie algebra.
  \item $\fb = \Biggl\{ \begin{pmatrix}
            * & * & \cdots & *\\
              & * & \cdots & *\\
              &   & \ddots & \vdots\\
              &   &        & *
          \end{pmatrix} \in \gl_n \Biggr\}$, the upper triangular matrices,
    form a solvable Lie algebra.

    \begin{exercise}[Essential]\leavevmode
    \begin{alist}
      \item Compute the central and derived series to check that these
        algebras are nilpotent and solvable respectively.
      \item Compute the centre of these algebras.
    \end{alist}
    \end{exercise}
  \item Let $W$ be a symplectic vector space with inner product
    $\langle\,,\,\rangle$ (recall that a symplectic vector space is equipped
    with a nondegenerate bilinear antisymmetric form, and there is essentially
    one example: given a finite-dimensional vector space $L$ let
    $W = L + L^*$ with inner product
    $\langle L, L\rangle = \langle L^*, L^*\rangle = 0$ and
    $\langle v, v^*\rangle = -\langle v^*, v\rangle = v^*(v)$). Define the
    \emph{Heisenberg Lie algebra} as
    \[
      \Heis_W = W \oplus \C c, \qquad [w, w'] = \langle w, w'\rangle c, \quad [c, w] = 0.
    \]

    \begin{exercise}
    Show this is a Lie algebra, and that it is nilpotent.
    \end{exercise}

    This is the most important nilpotent Lie algebra occurring in nature. For
    example, let $L = \C$; then $\Heis_W = \C p + \C q + \C c$ with
    $[p, q] = c$, $[p, c] = [q, c] = 0$ (as we know from classifying the
    $3$-dimensional Lie algebras!). Show that this has a representation on
    $\C[x]$ via $q \mapsto$ multiplication by $x$, $p \mapsto \partial/\partial x$,
    $c \mapsto 1$.

    (For a general vector space $L$, the basis vectors $v_1, \dots, v_n$ map
    to multiplication by $x_1, \dots, x_n$, and their duals map to
    $\partial/\partial x_1, \dots, \partial/\partial x_n$.)
\end{enumerate}
\end{example}

\begin{properties}\leavevmode
\begin{enumerate}
  \item Subalgebras and quotients of solvable Lie algebras are solvable.
    Subalgebras and quotients of nilpotent Lie algebras are nilpotent.
  \item Let $\g$ be a Lie algebra, and let $\h$ be an ideal of $\g$. Then $\g$
    is solvable if and only if $\h$ and $\g/\h$ are both solvable. (That is,
    solvable Lie algebras are built out of abelian Lie algebras, i.e.\ there
    is a refinement of the derived series such that the subquotients are
    $1$-dimensional and therefore abelian.)
  \item A nonzero $\g$ is nilpotent iff the centre of $\g$ is nonzero and
    $\g/(\text{centre of } \g)$ is nilpotent. Indeed, if $\g$ is nilpotent,
    then the central series is
    \[
      \g \supsetneq \g^1 \supsetneq \cdots \supsetneq \g^{n-1} \supsetneq \g^n = 0,
    \]
    and since $\g^n = [\g^{n-1}, \g]$ we must have $\g^{n-1}$ contained in the
    centre of $\g$.
  \item $\g$ is nilpotent iff $\ad \g \subseteq \gl(\g)$ is nilpotent, as we
    have an exact sequence
    \[
      0 \to \text{centre of } \g \hookrightarrow \g \twoheadrightarrow \ad \g \to 0.
    \]
\end{enumerate}
\end{properties}

\begin{exercise}
Prove the properties. (The proofs should all be easy and tedious.)
\end{exercise}

We now state a nontrivial result, which we will not prove.

\begin{theorem}[Lie's theorem]
Let $\g \subseteq \gl(V)$ be a solvable Lie algebra, where $V \neq 0$ is
finite-dimensional, and let the base field $k$ be algebraically closed and
have characteristic zero. Then there exists a basis $v_1, \dots, v_n$ of $V$
with respect to which $\g \subseteq \fb(V)$, i.e.\ the matrices of all elements
of $\g$ are upper triangular.

Equivalently, there exist a linear function $\lambda\colon \g \to k$ and a
nonzero $v \in V$ such that $xv = \lambda(x)v$ for all $x \in \g$, i.e.\ $v$ is
a common eigenvector for all of $\g$, i.e.\ $kv$ is a one-dimensional
subrepresentation.
\end{theorem}

In particular, the only irreducible finite-dimensional representations of a
solvable $\g$ are one-dimensional.

\begin{exercise}
Show that these are actually equivalent. (One direction should be obvious; in
the other, quotient by $v_1$ and repeat.)
\end{exercise}

\begin{exercise}
Show that we need $k$ to be algebraically closed.

Show that we need $k$ to have characteristic $0$. [\emph{Hint:} let $\g$ be
the $3$-dimensional Heisenberg Lie algebra, and show that $k[x]/x^p$ is a
finite-dimensional irreducible representation of dimension markedly bigger
than $1$.]
\end{exercise}

\begin{corollary}
In characteristic $0$, if $\g$ is a solvable finite-dimensional Lie algebra,
then $[\g, \g]$ is nilpotent.
\end{corollary}

\begin{exercise}
Find a counterexample in characteristic $p$.
\end{exercise}

\begin{proof}
Apply Lie's theorem to the adjoint representation $\g \to \End(\g)$. With
respect to some basis, $\ad \g \subseteq \fb(\g)$. Since $[\fb, \fb] \subseteq \fn$
(note the diagonal entries cancel out when we take the bracket), we see that
$[\ad \g, \ad \g]$ is nilpotent. Since $\ad[\g, \g] = [\ad \g, \ad \g]$ (it's a
representation!), we see that $\ad[\g, \g]$ is nilpotent, and therefore so is
$[\g, \g]$.
\end{proof}

\begin{theorem}[Engel's theorem]
Let the base field $k$ be arbitrary, and let $\g$ be finite-dimensional. Then
$\g$ is nilpotent iff $\ad \g$ consists of nilpotent endomorphisms of $\g$
(i.e., for all $x \in \g$, $\ad x$ is nilpotent).

Equivalently, if $V \neq 0$ is a finite-dimensional representation of $\g$ such
that all elements of $\g$ act on $V$ by nilpotent endomorphisms, then there
exists a nonzero $v \in V$ such that $x(v) = 0$ for all $x \in \g$, i.e.\ $V$
has the trivial representation as a subrepresentation.

Equivalently, in a suitable basis every $x \in \g$ acts by a strictly upper
triangular matrix.
\end{theorem}

\begin{exercise}
Show that these are actually equivalent.
\end{exercise}

\lecture{7}

\begin{definition}
A symmetric bilinear form $(\,,\,)\colon \g \times \g \to k$ is
\emph{invariant} if $([x, y], z) = (x, [y, z])$.
\end{definition}

\begin{exercise}
If $\fa \subseteq \g$ is an ideal and $(\,,\,)$ is an invariant bilinear form,
then $\fa^\perp$ is an ideal.
\end{exercise}

\begin{definition}
If $V$ is a finite-dimensional representation of $\g$, i.e.\ if
$\rho\colon \g \to \gl(V)$ is a homomorphism, define the \emph{trace form} to
be $(x, y)_V = \tr\bigl(\rho(x)\rho(y)\colon V \to V\bigr)$.
\end{definition}

\begin{exercise}
Check that the trace form is symmetric, bilinear, and invariant. (They should
all be obvious.)
\end{exercise}

\begin{example}
$(\,,\,)_{\ad}$ is the ``Killing form'' (named after a person, not after an
action). This is the trace form attached to the adjoint representation. That
is, $(x, y)_{\ad} = \tr(\ad x \ad y\colon \g \to \g)$.
\end{example}

\begin{theorem}[Cartan's criterion]
Let $\g \subseteq \gl(V)$ and let $\charac k = 0$. Then $\g$ is solvable if and
only if $(x, y)_V = 0$ for every $x \in \g$ and $y \in [\g, \g]$. (That is,
$[\g, \g] \subseteq \g^\perp$.)
\end{theorem}

\begin{exercise}
Lie's theorem gives us one direction. Indeed, if $\g$ is solvable then we can
take a basis in which $x$ is an upper triangular matrix, $y$ is strictly upper
triangular, and then $xy$ and $yx$ both have zeros on the diagonal (and thus
trace $0$). That is, if $\g$ is solvable and nonabelian, all trace forms are
degenerate. (The exercise is to convince yourself that this is true.)
\end{exercise}

\begin{corollary}
$\g$ is solvable iff $(\g, [\g, \g])_{\ad} = 0$.
\end{corollary}

\begin{proof}
If $\g$ is solvable, Lie's theorem gives the trace to be zero. Conversely,
Cartan's criterion tells us that $\ad \g = \g/(\text{centre of } \g)$ is
solvable, and therefore so is $\g$.
\end{proof}

\begin{exercise}
Not every invariant form is a trace form. Let
$\tilde\Heis = \C\langle p, q, c, d\rangle$ where $[c, \tilde\Heis] = 0$,
$[p, q] = c$, $[d, p] = p$, $[d, q] = -q$.
\begin{enumerate}
  \item Construct a nondegenerate invariant form on $\tilde\Heis$.
  \item Show $\tilde\Heis$ is solvable.
  \item Extend the representation of $\langle c, p, q\rangle = \Heis$ on
    $k[x]$ to a representation of $\tilde\Heis$.
\end{enumerate}
\end{exercise}

\begin{definition}
The \emph{radical} of $\g$, $R(\g)$, is the maximal solvable ideal in $\g$.
\end{definition}

\begin{exercise}\leavevmode
\begin{enumerate}
  \item Show $R(\g)$ is the sum of all solvable ideals in $\g$ (i.e., show
    that the sum of solvable ideals is solvable).
  \item Show $R(\g/R(\g)) = 0$.
\end{enumerate}
\end{exercise}

\begin{theorem}
In characteristic $0$, the following are equivalent:
\begin{enumerate}
  \item $\g$ is semisimple;
  \item $R(\g) = 0$;
  \item the Killing form $(\,,\,)_{\ad}$ is nondegenerate (the Killing
    criterion).
\end{enumerate}
Moreover, if $\g$ is semisimple, then every derivation $D\colon \g \to \g$ is
inner, but not conversely (where we define the terms ``derivation'' and
``inner'' below).
\end{theorem}

\begin{definition}
A \emph{derivation} $D\colon \g \to \g$ is a linear map satisfying
$D[x, y] = [Dx, y] + [x, Dy]$.
\end{definition}

\begin{example}
$\ad(x)$ is a derivation for any $x$. Derivations of the form $\ad(x)$ for some
$x \in \g$ are called \emph{inner}.
\end{example}

\begin{remark}
If $\g$ is any Lie algebra, we have the exact sequence
\[
  0 \to R(\g) \hookrightarrow \g \twoheadrightarrow \g/R(\g) \to 0,
\]
where $R(\g)$ is a solvable ideal, and $\g/R(\g)$ is semisimple. That is, any
Lie algebra has a maximal semisimple quotient, and the kernel is a solvable
ideal. This shows you how much nicer the theory of Lie algebras is than the
corresponding theory for finite groups. In particular, the corresponding
statement for finite groups is essentially equivalent to the classification of
the finite simple groups.
\end{remark}

A theorem that we will neither prove nor use, but which is pretty:

\begin{theorem}[Levi's theorem]
In characteristic $0$, a stronger result is true. Namely, the exact sequence
splits, i.e.\ there exists a subalgebra $\fs \subseteq \g$ such that
$\fs \cong \g/R(\g)$. (This subalgebra is not canonical; in particular, it is
not an ideal of $\g$.) That is, $\g$ can be written as
$\g = \fs \ltimes R(\g)$.
\end{theorem}

\begin{exercise}
Show that this fails in characteristic $p$. Let $\g = \sl_p(\Fpbar)$; show
$R(\g) = \Fpbar I$ (note that the identity matrix in dimension $p$ has trace
$0$), but that there is no complement to $R(\g)$ that is an algebra.
\end{exercise}

\begin{proof}[Proof of Theorem~4.7]
First, notice that $R(\g) = 0$ if and only if $\g$ has no nonzero abelian
ideals. (In one direction, an abelian ideal is solvable; in the other, notice
that the last nonzero term of the derived series of $R(\g)$ is an abelian
ideal of $\g$.)

From now on, $(\,,\,)$ refers to the Killing form $(\,,\,)_{\ad}$.

(3) $\implies$ (2): we show that if $\fa$ is an abelian ideal of $\g$, then
$\fa \subseteq \g^\perp$. (Since the Killing form is nondegenerate, this means
$\fa = 0$, and therefore $R(\g) = 0$.)

Write $\g = \fa + \h$ where $\h$ is a vector space complement to $\fa$ (not
necessarily an ideal). For $a \in \fa$, $\ad a$ has block matrix
$\bigl(\begin{smallmatrix} 0 & * \\ 0 & 0 \end{smallmatrix}\bigr)$. For
$x \in \g$, $\ad x$ has block matrix
$\bigl(\begin{smallmatrix} * & * \\ 0 & * \end{smallmatrix}\bigr)$. (This is
because $\fa$ is an ideal, so $[a, x] \in \fa$ for all $x \in \g$, and moreover
$\ad a$ acts by $0$ on $\fa$.) Therefore,
\[
  \tr(\ad a \cdot \ad x) = \tr \begin{pmatrix} 0 & * \\ 0 & 0 \end{pmatrix} = 0,
\]
i.e.\ $(\fa, \g)_{\ad} = 0$. Since the Killing form is nondegenerate, $\fa = 0$.

(2) $\implies$ (3): Let $\fr \subseteq \g^\perp$ be an ideal (e.g.\
$\fr = \g^\perp$), and suppose $\fr \neq 0$. Then $\ad \fr \subseteq \gl(\g)$,
and $(x, y)_{\ad} = 0$ for all $x, y \in \fr$. By Cartan's criterion,
$\ad \fr \cong \fr/(\fr \cap \text{centre of } \g)$ is solvable, so $\fr$ is
solvable, contradicting $R(\g) = 0$.

\begin{exercise}
Show that $R(\g) \supseteq \g^\perp \supseteq [R(\g), R(\g)]$.
\end{exercise}

(2), (3) $\implies$ (1): Let $(\,,\,)_{\ad}$ be nondegenerate, and let
$\fa \subseteq \g$ be a minimal nonzero ideal. We claim that
$(\,,\,)_{\ad}|_{\fa}$ is either $0$ or nondegenerate. Indeed, the kernel of
$(\,,\,)_{\ad}|_{\fa}$ is
$\setc{x \in \fa}{(x, \fa)_{\ad} = 0} = \fa \cap \fa^\perp$, which is an ideal!

Cartan's criterion implies that $\fa$ is solvable if $(\,,\,)_{\ad}|_{\fa}$ is
$0$. Since $R(\g) = 0$, we conclude that $(\,,\,)_{\ad}|_{\fa}$ is
nondegenerate.

Therefore, $\g = \fa \oplus \fa^\perp$, where $\fa$ is a minimal ideal, i.e.\
simple. (Note that $R(\g) = 0$, so $\fa$ cannot be abelian.)

But now ideals of $\fa^\perp$ are ideals of $\g$, so we can apply the same
argument to $\fa^\perp$, since $R(\g) = 0 \implies R(\fa^\perp) = 0$.
Therefore, $\g = \bigoplus \fa_i$, where the $\fa_i$ are simple Lie algebras.
\end{proof}

\begin{exercise}
Show that if $\g$ is semisimple, then $\g$ is the direct sum of minimal ideals
\emph{in a unique manner}. That is, if $\g = \bigoplus \fa_i$ where the
$\fa_i$ are minimal, and $\fb$ is a minimal ideal of $\g$, then in fact
$\fb = \fa_i$ for some $i$. (Consider $\fb \cap \fa_j$.) Conclude that
(1) $\implies$ (2).
\end{exercise}

\lecture{8} Finally, we show that all derivations on semisimple Lie algebras
are inner. Let $\g$ be a semisimple Lie algebra, and let $D\colon \g \to \g$ be
a derivation. Consider the linear functional $l\colon \g \to k$ via
$x \mapsto \tr(D \cdot \ad x)$. Since $\g$ is semisimple, the Killing form is a
nondegenerate inner product giving an isomorphism with the dual, so there
exists $y \in \g$ such that $l(x) = (y, x)_{\ad}$. Our task is to show that
$D = \ad y$, i.e.\ that $E := D - \ad y = 0$.

This is equivalent to showing that $Ex = 0$ for all $x \in \g$, or
equivalently that $(Ex, z)_{\ad} = 0$ for all $z \in \g$. Now,
\[
  \ad(Ex) = E \cdot \ad x - \ad x \cdot E = [E, \ad x]\colon \g \to \g,
\]
since $\ad(Ex)(z) = [Ex, z] = E[x, z] - [x, Ez]$. Therefore,
\begin{align*}
  (Ex, z)_{\ad} &= \tr_\g(\ad(Ex) \cdot \ad z) = \tr_\g([E, \ad x] \cdot \ad z)\\
    &= \tr_\g(E \cdot [\ad x, \ad z]) = \tr_\g(E \cdot \ad[x, z]).
\end{align*}
But by definition of $E$, for any $a \in \g$ (here $a = [x, z]$),
$\tr_\g(E \cdot \ad a) = \tr_\g(D \cdot \ad a) - (y, a)_{\ad} = l(a) - l(a) = 0$.
\qed

\begin{exercise}\leavevmode
\begin{enumerate}
  \item A nonzero nilpotent Lie algebra always has non-inner derivations.
  \item $\g = \langle a, b\rangle$ with $[a, b] = b$ only has inner
    derivations. Thus, this condition doesn't characterise semisimplicity.
\end{enumerate}
\end{exercise}

\begin{exercise}\leavevmode
\begin{enumerate}
  \item Let $\g$ be a simple Lie algebra, and $(\,,\,)_1$ and $(\,,\,)_2$ two
    non-degenerate invariant bilinear forms. Show that there exists $\lambda$
    such that $(\,,\,)_1 = \lambda(\,,\,)_2$.
  \item Let $\g = \sl_n(\C)$. We will shortly show that this is a simple Lie
    algebra. We have two nondegenerate bilinear forms: the Killing form
    $(\,,\,)_{\ad}$ and $(A, B) = \tr(AB)$ in the standard representation.
    Compute $\lambda$.
\end{enumerate}
\end{exercise}
